% Overleaf-ready Results and Discussion section.
% Recommended packages:
% \usepackage{booktabs}
% \usepackage{threeparttable}
% \usepackage{multirow}
% \usepackage{makecell}
% \usepackage{adjustbox}
% \usepackage{array}
% \usepackage{amsmath}
%
% Usage:
% % Overleaf-ready tables for the SCI manuscript.
% Recommended packages in the main preamble:
% \usepackage{booktabs}
% \usepackage{threeparttable}
% \usepackage{multirow}
% \usepackage{makecell}
% \usepackage{adjustbox}
% \usepackage{array}

\begin{table*}[t]
\centering
\caption{Benchmark configuration and compared methods.}
\label{tab:benchmark_setup}
\begin{threeparttable}
\begin{tabular}{p{2.8cm} p{11.8cm}}
\toprule
Component & Specification \\
\midrule
Main backbone & Local \texttt{Qwen2.5-7B-Instruct} executed on an RTX 3090 GPU \\
Cross-model validation backbone & Commercial \texttt{gpt-5.4-mini-2026-03-17} \\
Products & 10 \\
Audience segments & Students, young professionals, family users, and price-sensitive consumers \\
Tasks & 40 product-audience tasks \\
Primary baselines & \texttt{B0\_Template}, \texttt{B1\_SingleShot\_API}, \texttt{B2\_OpenSource\_Only}, \texttt{B3\_PromptEngineered\_AI}, and \texttt{Ours\_EAI\_CO} \\
Ablations & \texttt{Ours\_without\_diversity}, \texttt{Ours\_without\_factual\_penalty}, \texttt{Ours\_without\_iterative\_loop}, and \texttt{Ours\_without\_audience\_modeling} \\
Search depth & 2 rounds, 2 candidates per round, and 1 elite candidate retained for refinement \\
Reward weights (main benchmark) & relevance $0.20$, clarity $0.16$, aesthetic $0.14$, audience fit $0.16$, predicted engagement $0.16$, diversity $0.08$, and brand safety $0.10$ \\
Main evaluated outputs & 360 creatives for the local main experiment; 200 creatives for the commercial cross-model validation \\
\bottomrule
\end{tabular}
\end{threeparttable}
\end{table*}

\begin{table*}[t]
\centering
\caption{Main benchmark results on the local \texttt{Qwen2.5-7B-Instruct} backbone. Best values among the primary methods are shown in bold.}
\label{tab:qwen_main_results}
\begin{adjustbox}{max width=\textwidth}
\begin{tabular}{lccccccccc}
\toprule
Method & Reward & Relevance & Clarity & Audience Fit & Diversity & Pred.\ Eng. & Brand Safety & Penalty & Latency (ms) \\
\midrule
\texttt{Ours\_EAI\_CO} & \textbf{0.8051} & \textbf{0.9122} & 0.9130 & \textbf{0.4719} & 0.7693 & \textbf{0.8238} & \textbf{0.9945} & 0.0010 & 11114.23 \\
\texttt{B3\_PromptEngineered\_AI} & 0.7568 & 0.7920 & \textbf{0.9165} & 0.4394 & 0.8600 & 0.7849 & 0.9395 & 0.0110 & 2727.48 \\
\texttt{B1\_SingleShot\_API} & 0.7494 & 0.7920 & 0.9130 & 0.4367 & \textbf{0.8660} & 0.7695 & 0.9175 & 0.0130 & 2780.30 \\
\texttt{B2\_OpenSource\_Only} & 0.7445 & 0.8050 & 0.8955 & 0.4312 & 0.8631 & 0.7646 & 0.9340 & 0.0120 & 2707.25 \\
\texttt{B0\_Template} & 0.6856 & 0.5580 & 0.9200 & 0.4042 & 1.0000 & 0.6570 & 1.0000 & 0.0000 & 1.00 \\
\bottomrule
\end{tabular}
\end{adjustbox}
\end{table*}

\begin{table*}[t]
\centering
\caption{Ablation results on the local \texttt{Qwen2.5-7B-Instruct} benchmark. Negative reward gaps indicate degradation relative to the full EAI-CO framework.}
\label{tab:qwen_ablation}
\begin{adjustbox}{max width=\textwidth}
\begin{tabular}{lcccccc}
\toprule
Method & Reward & Reward Gap vs. Ours & Audience Fit & Audience-Fit Gap & Pred.\ Eng. & Brand Safety \\
\midrule
\texttt{Ours\_EAI\_CO} & \textbf{0.8051} & 0.0000 & \textbf{0.4719} & 0.0000 & \textbf{0.8238} & 0.9945 \\
\texttt{w/o factual penalty} & 0.8006 & -0.0045 & 0.4692 & -0.0027 & 0.8185 & 0.8130 \\
\texttt{w/o audience modeling} & 0.7833 & -0.0218 & 0.3960 & -0.0759 & 0.8100 & \textbf{1.0000} \\
\texttt{w/o iterative loop} & 0.7760 & -0.0291 & 0.4421 & -0.0298 & 0.7959 & 0.9670 \\
\texttt{w/o diversity} & 0.7349 & -0.0702 & 0.4556 & -0.0163 & 0.8149 & 0.9890 \\
\bottomrule
\end{tabular}
\end{adjustbox}
\end{table*}

\begin{table*}[t]
\centering
\caption{Audience-wise reward on the local \texttt{Qwen2.5-7B-Instruct} benchmark. Best values within each audience segment are shown in bold.}
\label{tab:qwen_audience}
\begin{adjustbox}{max width=\textwidth}
\begin{tabular}{lcccc}
\toprule
Method & Students & Young Professionals & Family Users & Price-sensitive Consumers \\
\midrule
\texttt{Ours\_EAI\_CO} & \textbf{0.8098} & \textbf{0.8050} & \textbf{0.8147} & \textbf{0.7908} \\
\texttt{B3\_PromptEngineered\_AI} & 0.7169 & 0.7554 & 0.7785 & 0.7764 \\
\texttt{B1\_SingleShot\_API} & 0.7648 & 0.7359 & 0.7389 & 0.7580 \\
\texttt{B2\_OpenSource\_Only} & 0.7229 & 0.7385 & 0.7663 & 0.7503 \\
\texttt{B0\_Template} & 0.6797 & 0.6875 & 0.6992 & 0.6758 \\
\bottomrule
\end{tabular}
\end{adjustbox}
\end{table*}

\begin{table*}[t]
\centering
\caption{Cross-model validation results on \texttt{gpt-5.4-mini-2026-03-17}. Best values among the primary methods are shown in bold.}
\label{tab:gpt_cross_model}
\begin{adjustbox}{max width=\textwidth}
\begin{tabular}{lccccccccc}
\toprule
Method & Reward & Relevance & Clarity & Audience Fit & Diversity & Pred.\ Eng. & Brand Safety & Penalty & Latency (ms) \\
\midrule
\texttt{Ours\_EAI\_CO} & \textbf{0.9957} & \textbf{1.0000} & \textbf{0.7390} & \textbf{0.7562} & 0.8041 & \textbf{0.8506} & \textbf{1.0000} & 0.0000 & 17128.73 \\
\texttt{B2\_OpenSource\_Only} & 0.9751 & 0.9870 & 0.7365 & 0.6452 & 0.8631 & 0.8145 & \textbf{1.0000} & 0.0000 & 6443.85 \\
\texttt{B1\_SingleShot\_API} & 0.9736 & 0.9870 & 0.6935 & 0.7183 & \textbf{0.8660} & 0.8211 & 0.9890 & 0.0110 & 6188.20 \\
\texttt{B3\_PromptEngineered\_AI} & 0.9729 & 0.9838 & 0.6925 & 0.6994 & 0.8600 & 0.8276 & 0.9890 & 0.0120 & 10776.03 \\
\texttt{B0\_Template} & 0.8396 & 0.5580 & 0.9200 & 0.4042 & 1.0000 & 0.6570 & \textbf{1.0000} & 0.0000 & 1.00 \\
\bottomrule
\end{tabular}
\end{adjustbox}
\end{table*}

\begin{table*}[t]
\centering
\caption{Statistical comparison of EAI-CO against primary baselines and runtime profile. Left: paired significance tests on the local \texttt{Qwen2.5-7B-Instruct} benchmark. Right: average runtime and model-call profile.}
\label{tab:significance_cost}
\begin{minipage}[t]{0.60\textwidth}
\centering
\begin{threeparttable}
\begin{tabular}{lccc}
\toprule
Baseline & $\Delta$ Reward & $p$-value & $\Delta$ Audience Fit \\
\midrule
\texttt{B0\_Template} & +0.1195 & $9.09\times10^{-13}$ & +0.0677 \\
\texttt{B1\_SingleShot\_API} & +0.0557 & $1.34\times10^{-6}$ & +0.0352 \\
\texttt{B2\_OpenSource\_Only} & +0.0606 & $5.18\times10^{-8}$ & +0.0406 \\
\texttt{B3\_PromptEngineered\_AI} & +0.0483 & $6.58\times10^{-5}$ & +0.0325 \\
\bottomrule
\end{tabular}
\end{threeparttable}
\end{minipage}
\hfill
\begin{minipage}[t]{0.36\textwidth}
\centering
\begin{threeparttable}
\begin{tabular}{lcc}
\toprule
Method & Calls & Mean Latency (ms) \\
\midrule
\texttt{Ours\_EAI\_CO} & 8.0 & 11114.23 \\
\texttt{B3\_PromptEngineered\_AI} & 4.0 & 2727.48 \\
\texttt{B1\_SingleShot\_API} & 4.0 & 2780.30 \\
\texttt{B2\_OpenSource\_Only} & 3.0 & 2707.25 \\
\texttt{B0\_Template} & 0.0 & 1.00 \\
\bottomrule
\end{tabular}
\end{threeparttable}
\end{minipage}
\end{table*}

% % Overleaf-ready Results and Discussion section.
% Recommended packages:
% \usepackage{booktabs}
% \usepackage{threeparttable}
% \usepackage{multirow}
% \usepackage{makecell}
% \usepackage{adjustbox}
% \usepackage{array}
% \usepackage{amsmath}
%
% Usage:
% % Overleaf-ready tables for the SCI manuscript.
% Recommended packages in the main preamble:
% \usepackage{booktabs}
% \usepackage{threeparttable}
% \usepackage{multirow}
% \usepackage{makecell}
% \usepackage{adjustbox}
% \usepackage{array}

\begin{table*}[t]
\centering
\caption{Benchmark configuration and compared methods.}
\label{tab:benchmark_setup}
\begin{threeparttable}
\begin{tabular}{p{2.8cm} p{11.8cm}}
\toprule
Component & Specification \\
\midrule
Main backbone & Local \texttt{Qwen2.5-7B-Instruct} executed on an RTX 3090 GPU \\
Cross-model validation backbone & Commercial \texttt{gpt-5.4-mini-2026-03-17} \\
Products & 10 \\
Audience segments & Students, young professionals, family users, and price-sensitive consumers \\
Tasks & 40 product-audience tasks \\
Primary baselines & \texttt{B0\_Template}, \texttt{B1\_SingleShot\_API}, \texttt{B2\_OpenSource\_Only}, \texttt{B3\_PromptEngineered\_AI}, and \texttt{Ours\_EAI\_CO} \\
Ablations & \texttt{Ours\_without\_diversity}, \texttt{Ours\_without\_factual\_penalty}, \texttt{Ours\_without\_iterative\_loop}, and \texttt{Ours\_without\_audience\_modeling} \\
Search depth & 2 rounds, 2 candidates per round, and 1 elite candidate retained for refinement \\
Reward weights (main benchmark) & relevance $0.20$, clarity $0.16$, aesthetic $0.14$, audience fit $0.16$, predicted engagement $0.16$, diversity $0.08$, and brand safety $0.10$ \\
Main evaluated outputs & 360 creatives for the local main experiment; 200 creatives for the commercial cross-model validation \\
\bottomrule
\end{tabular}
\end{threeparttable}
\end{table*}

\begin{table*}[t]
\centering
\caption{Main benchmark results on the local \texttt{Qwen2.5-7B-Instruct} backbone. Best values among the primary methods are shown in bold.}
\label{tab:qwen_main_results}
\begin{adjustbox}{max width=\textwidth}
\begin{tabular}{lccccccccc}
\toprule
Method & Reward & Relevance & Clarity & Audience Fit & Diversity & Pred.\ Eng. & Brand Safety & Penalty & Latency (ms) \\
\midrule
\texttt{Ours\_EAI\_CO} & \textbf{0.8051} & \textbf{0.9122} & 0.9130 & \textbf{0.4719} & 0.7693 & \textbf{0.8238} & \textbf{0.9945} & 0.0010 & 11114.23 \\
\texttt{B3\_PromptEngineered\_AI} & 0.7568 & 0.7920 & \textbf{0.9165} & 0.4394 & 0.8600 & 0.7849 & 0.9395 & 0.0110 & 2727.48 \\
\texttt{B1\_SingleShot\_API} & 0.7494 & 0.7920 & 0.9130 & 0.4367 & \textbf{0.8660} & 0.7695 & 0.9175 & 0.0130 & 2780.30 \\
\texttt{B2\_OpenSource\_Only} & 0.7445 & 0.8050 & 0.8955 & 0.4312 & 0.8631 & 0.7646 & 0.9340 & 0.0120 & 2707.25 \\
\texttt{B0\_Template} & 0.6856 & 0.5580 & 0.9200 & 0.4042 & 1.0000 & 0.6570 & 1.0000 & 0.0000 & 1.00 \\
\bottomrule
\end{tabular}
\end{adjustbox}
\end{table*}

\begin{table*}[t]
\centering
\caption{Ablation results on the local \texttt{Qwen2.5-7B-Instruct} benchmark. Negative reward gaps indicate degradation relative to the full EAI-CO framework.}
\label{tab:qwen_ablation}
\begin{adjustbox}{max width=\textwidth}
\begin{tabular}{lcccccc}
\toprule
Method & Reward & Reward Gap vs. Ours & Audience Fit & Audience-Fit Gap & Pred.\ Eng. & Brand Safety \\
\midrule
\texttt{Ours\_EAI\_CO} & \textbf{0.8051} & 0.0000 & \textbf{0.4719} & 0.0000 & \textbf{0.8238} & 0.9945 \\
\texttt{w/o factual penalty} & 0.8006 & -0.0045 & 0.4692 & -0.0027 & 0.8185 & 0.8130 \\
\texttt{w/o audience modeling} & 0.7833 & -0.0218 & 0.3960 & -0.0759 & 0.8100 & \textbf{1.0000} \\
\texttt{w/o iterative loop} & 0.7760 & -0.0291 & 0.4421 & -0.0298 & 0.7959 & 0.9670 \\
\texttt{w/o diversity} & 0.7349 & -0.0702 & 0.4556 & -0.0163 & 0.8149 & 0.9890 \\
\bottomrule
\end{tabular}
\end{adjustbox}
\end{table*}

\begin{table*}[t]
\centering
\caption{Audience-wise reward on the local \texttt{Qwen2.5-7B-Instruct} benchmark. Best values within each audience segment are shown in bold.}
\label{tab:qwen_audience}
\begin{adjustbox}{max width=\textwidth}
\begin{tabular}{lcccc}
\toprule
Method & Students & Young Professionals & Family Users & Price-sensitive Consumers \\
\midrule
\texttt{Ours\_EAI\_CO} & \textbf{0.8098} & \textbf{0.8050} & \textbf{0.8147} & \textbf{0.7908} \\
\texttt{B3\_PromptEngineered\_AI} & 0.7169 & 0.7554 & 0.7785 & 0.7764 \\
\texttt{B1\_SingleShot\_API} & 0.7648 & 0.7359 & 0.7389 & 0.7580 \\
\texttt{B2\_OpenSource\_Only} & 0.7229 & 0.7385 & 0.7663 & 0.7503 \\
\texttt{B0\_Template} & 0.6797 & 0.6875 & 0.6992 & 0.6758 \\
\bottomrule
\end{tabular}
\end{adjustbox}
\end{table*}

\begin{table*}[t]
\centering
\caption{Cross-model validation results on \texttt{gpt-5.4-mini-2026-03-17}. Best values among the primary methods are shown in bold.}
\label{tab:gpt_cross_model}
\begin{adjustbox}{max width=\textwidth}
\begin{tabular}{lccccccccc}
\toprule
Method & Reward & Relevance & Clarity & Audience Fit & Diversity & Pred.\ Eng. & Brand Safety & Penalty & Latency (ms) \\
\midrule
\texttt{Ours\_EAI\_CO} & \textbf{0.9957} & \textbf{1.0000} & \textbf{0.7390} & \textbf{0.7562} & 0.8041 & \textbf{0.8506} & \textbf{1.0000} & 0.0000 & 17128.73 \\
\texttt{B2\_OpenSource\_Only} & 0.9751 & 0.9870 & 0.7365 & 0.6452 & 0.8631 & 0.8145 & \textbf{1.0000} & 0.0000 & 6443.85 \\
\texttt{B1\_SingleShot\_API} & 0.9736 & 0.9870 & 0.6935 & 0.7183 & \textbf{0.8660} & 0.8211 & 0.9890 & 0.0110 & 6188.20 \\
\texttt{B3\_PromptEngineered\_AI} & 0.9729 & 0.9838 & 0.6925 & 0.6994 & 0.8600 & 0.8276 & 0.9890 & 0.0120 & 10776.03 \\
\texttt{B0\_Template} & 0.8396 & 0.5580 & 0.9200 & 0.4042 & 1.0000 & 0.6570 & \textbf{1.0000} & 0.0000 & 1.00 \\
\bottomrule
\end{tabular}
\end{adjustbox}
\end{table*}

\begin{table*}[t]
\centering
\caption{Statistical comparison of EAI-CO against primary baselines and runtime profile. Left: paired significance tests on the local \texttt{Qwen2.5-7B-Instruct} benchmark. Right: average runtime and model-call profile.}
\label{tab:significance_cost}
\begin{minipage}[t]{0.60\textwidth}
\centering
\begin{threeparttable}
\begin{tabular}{lccc}
\toprule
Baseline & $\Delta$ Reward & $p$-value & $\Delta$ Audience Fit \\
\midrule
\texttt{B0\_Template} & +0.1195 & $9.09\times10^{-13}$ & +0.0677 \\
\texttt{B1\_SingleShot\_API} & +0.0557 & $1.34\times10^{-6}$ & +0.0352 \\
\texttt{B2\_OpenSource\_Only} & +0.0606 & $5.18\times10^{-8}$ & +0.0406 \\
\texttt{B3\_PromptEngineered\_AI} & +0.0483 & $6.58\times10^{-5}$ & +0.0325 \\
\bottomrule
\end{tabular}
\end{threeparttable}
\end{minipage}
\hfill
\begin{minipage}[t]{0.36\textwidth}
\centering
\begin{threeparttable}
\begin{tabular}{lcc}
\toprule
Method & Calls & Mean Latency (ms) \\
\midrule
\texttt{Ours\_EAI\_CO} & 8.0 & 11114.23 \\
\texttt{B3\_PromptEngineered\_AI} & 4.0 & 2727.48 \\
\texttt{B1\_SingleShot\_API} & 4.0 & 2780.30 \\
\texttt{B2\_OpenSource\_Only} & 3.0 & 2707.25 \\
\texttt{B0\_Template} & 0.0 & 1.00 \\
\bottomrule
\end{tabular}
\end{threeparttable}
\end{minipage}
\end{table*}

% % Overleaf-ready Results and Discussion section.
% Recommended packages:
% \usepackage{booktabs}
% \usepackage{threeparttable}
% \usepackage{multirow}
% \usepackage{makecell}
% \usepackage{adjustbox}
% \usepackage{array}
% \usepackage{amsmath}
%
% Usage:
% % Overleaf-ready tables for the SCI manuscript.
% Recommended packages in the main preamble:
% \usepackage{booktabs}
% \usepackage{threeparttable}
% \usepackage{multirow}
% \usepackage{makecell}
% \usepackage{adjustbox}
% \usepackage{array}

\begin{table*}[t]
\centering
\caption{Benchmark configuration and compared methods.}
\label{tab:benchmark_setup}
\begin{threeparttable}
\begin{tabular}{p{2.8cm} p{11.8cm}}
\toprule
Component & Specification \\
\midrule
Main backbone & Local \texttt{Qwen2.5-7B-Instruct} executed on an RTX 3090 GPU \\
Cross-model validation backbone & Commercial \texttt{gpt-5.4-mini-2026-03-17} \\
Products & 10 \\
Audience segments & Students, young professionals, family users, and price-sensitive consumers \\
Tasks & 40 product-audience tasks \\
Primary baselines & \texttt{B0\_Template}, \texttt{B1\_SingleShot\_API}, \texttt{B2\_OpenSource\_Only}, \texttt{B3\_PromptEngineered\_AI}, and \texttt{Ours\_EAI\_CO} \\
Ablations & \texttt{Ours\_without\_diversity}, \texttt{Ours\_without\_factual\_penalty}, \texttt{Ours\_without\_iterative\_loop}, and \texttt{Ours\_without\_audience\_modeling} \\
Search depth & 2 rounds, 2 candidates per round, and 1 elite candidate retained for refinement \\
Reward weights (main benchmark) & relevance $0.20$, clarity $0.16$, aesthetic $0.14$, audience fit $0.16$, predicted engagement $0.16$, diversity $0.08$, and brand safety $0.10$ \\
Main evaluated outputs & 360 creatives for the local main experiment; 200 creatives for the commercial cross-model validation \\
\bottomrule
\end{tabular}
\end{threeparttable}
\end{table*}

\begin{table*}[t]
\centering
\caption{Main benchmark results on the local \texttt{Qwen2.5-7B-Instruct} backbone. Best values among the primary methods are shown in bold.}
\label{tab:qwen_main_results}
\begin{adjustbox}{max width=\textwidth}
\begin{tabular}{lccccccccc}
\toprule
Method & Reward & Relevance & Clarity & Audience Fit & Diversity & Pred.\ Eng. & Brand Safety & Penalty & Latency (ms) \\
\midrule
\texttt{Ours\_EAI\_CO} & \textbf{0.8051} & \textbf{0.9122} & 0.9130 & \textbf{0.4719} & 0.7693 & \textbf{0.8238} & \textbf{0.9945} & 0.0010 & 11114.23 \\
\texttt{B3\_PromptEngineered\_AI} & 0.7568 & 0.7920 & \textbf{0.9165} & 0.4394 & 0.8600 & 0.7849 & 0.9395 & 0.0110 & 2727.48 \\
\texttt{B1\_SingleShot\_API} & 0.7494 & 0.7920 & 0.9130 & 0.4367 & \textbf{0.8660} & 0.7695 & 0.9175 & 0.0130 & 2780.30 \\
\texttt{B2\_OpenSource\_Only} & 0.7445 & 0.8050 & 0.8955 & 0.4312 & 0.8631 & 0.7646 & 0.9340 & 0.0120 & 2707.25 \\
\texttt{B0\_Template} & 0.6856 & 0.5580 & 0.9200 & 0.4042 & 1.0000 & 0.6570 & 1.0000 & 0.0000 & 1.00 \\
\bottomrule
\end{tabular}
\end{adjustbox}
\end{table*}

\begin{table*}[t]
\centering
\caption{Ablation results on the local \texttt{Qwen2.5-7B-Instruct} benchmark. Negative reward gaps indicate degradation relative to the full EAI-CO framework.}
\label{tab:qwen_ablation}
\begin{adjustbox}{max width=\textwidth}
\begin{tabular}{lcccccc}
\toprule
Method & Reward & Reward Gap vs. Ours & Audience Fit & Audience-Fit Gap & Pred.\ Eng. & Brand Safety \\
\midrule
\texttt{Ours\_EAI\_CO} & \textbf{0.8051} & 0.0000 & \textbf{0.4719} & 0.0000 & \textbf{0.8238} & 0.9945 \\
\texttt{w/o factual penalty} & 0.8006 & -0.0045 & 0.4692 & -0.0027 & 0.8185 & 0.8130 \\
\texttt{w/o audience modeling} & 0.7833 & -0.0218 & 0.3960 & -0.0759 & 0.8100 & \textbf{1.0000} \\
\texttt{w/o iterative loop} & 0.7760 & -0.0291 & 0.4421 & -0.0298 & 0.7959 & 0.9670 \\
\texttt{w/o diversity} & 0.7349 & -0.0702 & 0.4556 & -0.0163 & 0.8149 & 0.9890 \\
\bottomrule
\end{tabular}
\end{adjustbox}
\end{table*}

\begin{table*}[t]
\centering
\caption{Audience-wise reward on the local \texttt{Qwen2.5-7B-Instruct} benchmark. Best values within each audience segment are shown in bold.}
\label{tab:qwen_audience}
\begin{adjustbox}{max width=\textwidth}
\begin{tabular}{lcccc}
\toprule
Method & Students & Young Professionals & Family Users & Price-sensitive Consumers \\
\midrule
\texttt{Ours\_EAI\_CO} & \textbf{0.8098} & \textbf{0.8050} & \textbf{0.8147} & \textbf{0.7908} \\
\texttt{B3\_PromptEngineered\_AI} & 0.7169 & 0.7554 & 0.7785 & 0.7764 \\
\texttt{B1\_SingleShot\_API} & 0.7648 & 0.7359 & 0.7389 & 0.7580 \\
\texttt{B2\_OpenSource\_Only} & 0.7229 & 0.7385 & 0.7663 & 0.7503 \\
\texttt{B0\_Template} & 0.6797 & 0.6875 & 0.6992 & 0.6758 \\
\bottomrule
\end{tabular}
\end{adjustbox}
\end{table*}

\begin{table*}[t]
\centering
\caption{Cross-model validation results on \texttt{gpt-5.4-mini-2026-03-17}. Best values among the primary methods are shown in bold.}
\label{tab:gpt_cross_model}
\begin{adjustbox}{max width=\textwidth}
\begin{tabular}{lccccccccc}
\toprule
Method & Reward & Relevance & Clarity & Audience Fit & Diversity & Pred.\ Eng. & Brand Safety & Penalty & Latency (ms) \\
\midrule
\texttt{Ours\_EAI\_CO} & \textbf{0.9957} & \textbf{1.0000} & \textbf{0.7390} & \textbf{0.7562} & 0.8041 & \textbf{0.8506} & \textbf{1.0000} & 0.0000 & 17128.73 \\
\texttt{B2\_OpenSource\_Only} & 0.9751 & 0.9870 & 0.7365 & 0.6452 & 0.8631 & 0.8145 & \textbf{1.0000} & 0.0000 & 6443.85 \\
\texttt{B1\_SingleShot\_API} & 0.9736 & 0.9870 & 0.6935 & 0.7183 & \textbf{0.8660} & 0.8211 & 0.9890 & 0.0110 & 6188.20 \\
\texttt{B3\_PromptEngineered\_AI} & 0.9729 & 0.9838 & 0.6925 & 0.6994 & 0.8600 & 0.8276 & 0.9890 & 0.0120 & 10776.03 \\
\texttt{B0\_Template} & 0.8396 & 0.5580 & 0.9200 & 0.4042 & 1.0000 & 0.6570 & \textbf{1.0000} & 0.0000 & 1.00 \\
\bottomrule
\end{tabular}
\end{adjustbox}
\end{table*}

\begin{table*}[t]
\centering
\caption{Statistical comparison of EAI-CO against primary baselines and runtime profile. Left: paired significance tests on the local \texttt{Qwen2.5-7B-Instruct} benchmark. Right: average runtime and model-call profile.}
\label{tab:significance_cost}
\begin{minipage}[t]{0.60\textwidth}
\centering
\begin{threeparttable}
\begin{tabular}{lccc}
\toprule
Baseline & $\Delta$ Reward & $p$-value & $\Delta$ Audience Fit \\
\midrule
\texttt{B0\_Template} & +0.1195 & $9.09\times10^{-13}$ & +0.0677 \\
\texttt{B1\_SingleShot\_API} & +0.0557 & $1.34\times10^{-6}$ & +0.0352 \\
\texttt{B2\_OpenSource\_Only} & +0.0606 & $5.18\times10^{-8}$ & +0.0406 \\
\texttt{B3\_PromptEngineered\_AI} & +0.0483 & $6.58\times10^{-5}$ & +0.0325 \\
\bottomrule
\end{tabular}
\end{threeparttable}
\end{minipage}
\hfill
\begin{minipage}[t]{0.36\textwidth}
\centering
\begin{threeparttable}
\begin{tabular}{lcc}
\toprule
Method & Calls & Mean Latency (ms) \\
\midrule
\texttt{Ours\_EAI\_CO} & 8.0 & 11114.23 \\
\texttt{B3\_PromptEngineered\_AI} & 4.0 & 2727.48 \\
\texttt{B1\_SingleShot\_API} & 4.0 & 2780.30 \\
\texttt{B2\_OpenSource\_Only} & 3.0 & 2707.25 \\
\texttt{B0\_Template} & 0.0 & 1.00 \\
\bottomrule
\end{tabular}
\end{threeparttable}
\end{minipage}
\end{table*}

% % Overleaf-ready Results and Discussion section.
% Recommended packages:
% \usepackage{booktabs}
% \usepackage{threeparttable}
% \usepackage{multirow}
% \usepackage{makecell}
% \usepackage{adjustbox}
% \usepackage{array}
% \usepackage{amsmath}
%
% Usage:
% \input{manuscript/overleaf_tables}
% \input{manuscript/overleaf_results_discussion}

\section{Results}

\subsection{Main Benchmark Results on the Local Backbone}

Table~\ref{tab:qwen_main_results} reports the main benchmark results obtained on the local \texttt{Qwen2.5-7B-Instruct} backbone. EAI-CO achieved the highest overall reward (\textbf{0.8051}), outperforming \texttt{B3\_PromptEngineered\_AI} (0.7568), \texttt{B1\_SingleShot\_API} (0.7494), \texttt{B2\_OpenSource\_Only} (0.7445), and \texttt{B0\_Template} (0.6856). As summarized in Table~\ref{tab:significance_cost}, the reward gains over all primary baselines were statistically significant. Relative to \texttt{B0\_Template}, EAI-CO improved reward by +0.1195 ($p=9.09\times10^{-13}$). The gains over the stronger single-pass baselines also remained significant: +0.0557 over \texttt{B1\_SingleShot\_API} ($p=1.34\times10^{-6}$), +0.0606 over \texttt{B2\_OpenSource\_Only} ($p=5.18\times10^{-8}$), and +0.0483 over \texttt{B3\_PromptEngineered\_AI} ($p=6.58\times10^{-5}$).

Beyond the scalar reward, the full framework achieved the best audience-fit score (0.4719) and the strongest brand-safety score (0.9945) among the primary methods. This indicates that the observed advantage is not restricted to a single dimension such as relevance or engagement, but reflects a more balanced optimization outcome.

\subsection{Ablation Analysis}

Table~\ref{tab:qwen_ablation} presents the ablation study on the same benchmark. The full method outperformed all ablated variants. The largest degradation occurred when diversity preservation was removed (-0.0702), indicating that search-space breadth is a major contributor to the final gain. Removing the iterative loop reduced reward by -0.0291, showing that controlled refinement contributes beyond one-pass generation. Removing audience modeling lowered reward by -0.0218 and reduced audience fit from 0.4719 to 0.3960, directly supporting the claim that persona conditioning is functionally important.

The factual-penalty ablation remained close to the full model (-0.0045). This result should be interpreted carefully. It does not imply that compliance-aware controls are unnecessary. Rather, it indicates that factual safeguards can be integrated with only a small tradeoff in persuasive reward while preserving a more disciplined search process.

\subsection{Audience-wise Robustness}

Table~\ref{tab:qwen_audience} reports the audience-wise reward values on the local benchmark. EAI-CO ranked first in all four audience groups: family users (0.8147), students (0.8098), young professionals (0.8050), and price-sensitive consumers (0.7908). This consistency reduces the concern that the overall effect is driven by a single favorable subgroup and suggests that the exploratory optimization strategy remains effective across materially different audience profiles.

\subsection{Runtime Profile}

The runtime profile is summarized in Table~\ref{tab:significance_cost}. EAI-CO required 8.0 mean model calls per task and 11114.23 ms mean latency, whereas the single-pass baselines required approximately 3--4 model calls and 2700--2800 ms mean latency. This overhead is expected because the framework explicitly trades additional inference budget for exploratory search quality. For the intended offline campaign-design setting, this tradeoff remains operationally acceptable.

\subsection{Cross-model Validation on \texttt{gpt-5.4-mini-2026-03-17}}

Table~\ref{tab:gpt_cross_model} presents the fixed-subset cross-model validation on \texttt{gpt-5.4-mini-2026-03-17}. The same ranking pattern observed on the local backbone was preserved: EAI-CO achieved the highest reward (\textbf{0.9957}), followed by \texttt{B2\_OpenSource\_Only} (0.9751), \texttt{B1\_SingleShot\_API} (0.9736), \texttt{B3\_PromptEngineered\_AI} (0.9729), and \texttt{B0\_Template} (0.8396). The gains over all baselines remained statistically significant, with reward improvements of +0.1561 ($p=1.54\times10^{-8}$) over \texttt{B0\_Template}, +0.0221 ($p=1.41\times10^{-3}$) over \texttt{B1\_SingleShot\_API}, +0.0206 ($p=6.42\times10^{-4}$) over \texttt{B2\_OpenSource\_Only}, and +0.0228 ($p=1.01\times10^{-3}$) over \texttt{B3\_PromptEngineered\_AI}.

The audience-wise rankings were also preserved, with EAI-CO remaining first for family users (0.9950), price-sensitive consumers (0.9938), students (1.0000), and young professionals (0.9939). These results support the claim that the advantage of the proposed exploratory optimization framework transfers to a stronger commercial model under the same task protocol.

\section{Discussion}

\subsection{The Main Advantage Comes from Optimization Structure}

The most stable pattern across the local benchmark and the commercial-model validation is that exploratory optimization outperforms one-shot generation. This result suggests that the central contribution of EAI-CO lies less in any specific generator and more in the optimization structure itself. In digital-media advertising, creative quality should therefore be treated as a search-and-selection problem rather than as a direct text-generation problem.

\subsection{Why Diversity and Iteration Matter}

The ablation study shows that diversity preservation and iterative refinement are the main structural sources of improvement. This is consistent with how advertising work is performed in practice. Strong campaign creatives rarely emerge from the first acceptable draft; they emerge from comparing strategically different variants and refining the strongest direction. The present benchmark makes this process measurable.

\subsection{Audience Modeling Has Measurable Functional Value}

The audience-modeling ablation provides direct empirical evidence that personalization is not simply a rhetorical feature. Once persona conditioning is removed, both reward and audience fit decline. This implies that the framework is not merely generating generic persuasive copy but is measurably exploiting audience-specific information to alter the optimization trajectory.

\subsection{Compliance-aware Scoring Is a Practical Requirement}

The factual-penalty ablation remained close to the full framework, which indicates that compliance-aware optimization can be introduced without substantially damaging persuasive performance. This is a practically desirable result. In real digital-media settings, a method that produces slightly stronger promotional language at the cost of weaker factual discipline would not necessarily be preferable.

\subsection{Cross-model Transfer Strengthens the Paper}

The commercial-model experiment should be interpreted as a transferability test rather than as the main source of effect-size contrast. Because \texttt{gpt-5.4-mini} is a stronger generator under the current evaluator, reward values become compressed near the upper bound. Nevertheless, the preservation of the full ranking pattern materially strengthens the paper by reducing the risk that the main result is merely an artifact of one local open-source backbone.

\subsection{Limitations}

The present study does not include real platform CTR or conversion outcomes, and the current submission package does not yet include a completed human-evaluation study. Accordingly, the paper should restrict its claims to automatic reward superiority, predicted engagement, audience-conditioned quality, and cross-model ranking consistency. In addition, the current benchmark remains smaller than a production-scale industrial evaluation environment, and future work should expand both product diversity and deployment realism.


\section{Results}

\subsection{Main Benchmark Results on the Local Backbone}

Table~\ref{tab:qwen_main_results} reports the main benchmark results obtained on the local \texttt{Qwen2.5-7B-Instruct} backbone. EAI-CO achieved the highest overall reward (\textbf{0.8051}), outperforming \texttt{B3\_PromptEngineered\_AI} (0.7568), \texttt{B1\_SingleShot\_API} (0.7494), \texttt{B2\_OpenSource\_Only} (0.7445), and \texttt{B0\_Template} (0.6856). As summarized in Table~\ref{tab:significance_cost}, the reward gains over all primary baselines were statistically significant. Relative to \texttt{B0\_Template}, EAI-CO improved reward by +0.1195 ($p=9.09\times10^{-13}$). The gains over the stronger single-pass baselines also remained significant: +0.0557 over \texttt{B1\_SingleShot\_API} ($p=1.34\times10^{-6}$), +0.0606 over \texttt{B2\_OpenSource\_Only} ($p=5.18\times10^{-8}$), and +0.0483 over \texttt{B3\_PromptEngineered\_AI} ($p=6.58\times10^{-5}$).

Beyond the scalar reward, the full framework achieved the best audience-fit score (0.4719) and the strongest brand-safety score (0.9945) among the primary methods. This indicates that the observed advantage is not restricted to a single dimension such as relevance or engagement, but reflects a more balanced optimization outcome.

\subsection{Ablation Analysis}

Table~\ref{tab:qwen_ablation} presents the ablation study on the same benchmark. The full method outperformed all ablated variants. The largest degradation occurred when diversity preservation was removed (-0.0702), indicating that search-space breadth is a major contributor to the final gain. Removing the iterative loop reduced reward by -0.0291, showing that controlled refinement contributes beyond one-pass generation. Removing audience modeling lowered reward by -0.0218 and reduced audience fit from 0.4719 to 0.3960, directly supporting the claim that persona conditioning is functionally important.

The factual-penalty ablation remained close to the full model (-0.0045). This result should be interpreted carefully. It does not imply that compliance-aware controls are unnecessary. Rather, it indicates that factual safeguards can be integrated with only a small tradeoff in persuasive reward while preserving a more disciplined search process.

\subsection{Audience-wise Robustness}

Table~\ref{tab:qwen_audience} reports the audience-wise reward values on the local benchmark. EAI-CO ranked first in all four audience groups: family users (0.8147), students (0.8098), young professionals (0.8050), and price-sensitive consumers (0.7908). This consistency reduces the concern that the overall effect is driven by a single favorable subgroup and suggests that the exploratory optimization strategy remains effective across materially different audience profiles.

\subsection{Runtime Profile}

The runtime profile is summarized in Table~\ref{tab:significance_cost}. EAI-CO required 8.0 mean model calls per task and 11114.23 ms mean latency, whereas the single-pass baselines required approximately 3--4 model calls and 2700--2800 ms mean latency. This overhead is expected because the framework explicitly trades additional inference budget for exploratory search quality. For the intended offline campaign-design setting, this tradeoff remains operationally acceptable.

\subsection{Cross-model Validation on \texttt{gpt-5.4-mini-2026-03-17}}

Table~\ref{tab:gpt_cross_model} presents the fixed-subset cross-model validation on \texttt{gpt-5.4-mini-2026-03-17}. The same ranking pattern observed on the local backbone was preserved: EAI-CO achieved the highest reward (\textbf{0.9957}), followed by \texttt{B2\_OpenSource\_Only} (0.9751), \texttt{B1\_SingleShot\_API} (0.9736), \texttt{B3\_PromptEngineered\_AI} (0.9729), and \texttt{B0\_Template} (0.8396). The gains over all baselines remained statistically significant, with reward improvements of +0.1561 ($p=1.54\times10^{-8}$) over \texttt{B0\_Template}, +0.0221 ($p=1.41\times10^{-3}$) over \texttt{B1\_SingleShot\_API}, +0.0206 ($p=6.42\times10^{-4}$) over \texttt{B2\_OpenSource\_Only}, and +0.0228 ($p=1.01\times10^{-3}$) over \texttt{B3\_PromptEngineered\_AI}.

The audience-wise rankings were also preserved, with EAI-CO remaining first for family users (0.9950), price-sensitive consumers (0.9938), students (1.0000), and young professionals (0.9939). These results support the claim that the advantage of the proposed exploratory optimization framework transfers to a stronger commercial model under the same task protocol.

\section{Discussion}

\subsection{The Main Advantage Comes from Optimization Structure}

The most stable pattern across the local benchmark and the commercial-model validation is that exploratory optimization outperforms one-shot generation. This result suggests that the central contribution of EAI-CO lies less in any specific generator and more in the optimization structure itself. In digital-media advertising, creative quality should therefore be treated as a search-and-selection problem rather than as a direct text-generation problem.

\subsection{Why Diversity and Iteration Matter}

The ablation study shows that diversity preservation and iterative refinement are the main structural sources of improvement. This is consistent with how advertising work is performed in practice. Strong campaign creatives rarely emerge from the first acceptable draft; they emerge from comparing strategically different variants and refining the strongest direction. The present benchmark makes this process measurable.

\subsection{Audience Modeling Has Measurable Functional Value}

The audience-modeling ablation provides direct empirical evidence that personalization is not simply a rhetorical feature. Once persona conditioning is removed, both reward and audience fit decline. This implies that the framework is not merely generating generic persuasive copy but is measurably exploiting audience-specific information to alter the optimization trajectory.

\subsection{Compliance-aware Scoring Is a Practical Requirement}

The factual-penalty ablation remained close to the full framework, which indicates that compliance-aware optimization can be introduced without substantially damaging persuasive performance. This is a practically desirable result. In real digital-media settings, a method that produces slightly stronger promotional language at the cost of weaker factual discipline would not necessarily be preferable.

\subsection{Cross-model Transfer Strengthens the Paper}

The commercial-model experiment should be interpreted as a transferability test rather than as the main source of effect-size contrast. Because \texttt{gpt-5.4-mini} is a stronger generator under the current evaluator, reward values become compressed near the upper bound. Nevertheless, the preservation of the full ranking pattern materially strengthens the paper by reducing the risk that the main result is merely an artifact of one local open-source backbone.

\subsection{Limitations}

The present study does not include real platform CTR or conversion outcomes, and the current submission package does not yet include a completed human-evaluation study. Accordingly, the paper should restrict its claims to automatic reward superiority, predicted engagement, audience-conditioned quality, and cross-model ranking consistency. In addition, the current benchmark remains smaller than a production-scale industrial evaluation environment, and future work should expand both product diversity and deployment realism.


\section{Results}

\subsection{Main Benchmark Results on the Local Backbone}

Table~\ref{tab:qwen_main_results} reports the main benchmark results obtained on the local \texttt{Qwen2.5-7B-Instruct} backbone. EAI-CO achieved the highest overall reward (\textbf{0.8051}), outperforming \texttt{B3\_PromptEngineered\_AI} (0.7568), \texttt{B1\_SingleShot\_API} (0.7494), \texttt{B2\_OpenSource\_Only} (0.7445), and \texttt{B0\_Template} (0.6856). As summarized in Table~\ref{tab:significance_cost}, the reward gains over all primary baselines were statistically significant. Relative to \texttt{B0\_Template}, EAI-CO improved reward by +0.1195 ($p=9.09\times10^{-13}$). The gains over the stronger single-pass baselines also remained significant: +0.0557 over \texttt{B1\_SingleShot\_API} ($p=1.34\times10^{-6}$), +0.0606 over \texttt{B2\_OpenSource\_Only} ($p=5.18\times10^{-8}$), and +0.0483 over \texttt{B3\_PromptEngineered\_AI} ($p=6.58\times10^{-5}$).

Beyond the scalar reward, the full framework achieved the best audience-fit score (0.4719) and the strongest brand-safety score (0.9945) among the primary methods. This indicates that the observed advantage is not restricted to a single dimension such as relevance or engagement, but reflects a more balanced optimization outcome.

\subsection{Ablation Analysis}

Table~\ref{tab:qwen_ablation} presents the ablation study on the same benchmark. The full method outperformed all ablated variants. The largest degradation occurred when diversity preservation was removed (-0.0702), indicating that search-space breadth is a major contributor to the final gain. Removing the iterative loop reduced reward by -0.0291, showing that controlled refinement contributes beyond one-pass generation. Removing audience modeling lowered reward by -0.0218 and reduced audience fit from 0.4719 to 0.3960, directly supporting the claim that persona conditioning is functionally important.

The factual-penalty ablation remained close to the full model (-0.0045). This result should be interpreted carefully. It does not imply that compliance-aware controls are unnecessary. Rather, it indicates that factual safeguards can be integrated with only a small tradeoff in persuasive reward while preserving a more disciplined search process.

\subsection{Audience-wise Robustness}

Table~\ref{tab:qwen_audience} reports the audience-wise reward values on the local benchmark. EAI-CO ranked first in all four audience groups: family users (0.8147), students (0.8098), young professionals (0.8050), and price-sensitive consumers (0.7908). This consistency reduces the concern that the overall effect is driven by a single favorable subgroup and suggests that the exploratory optimization strategy remains effective across materially different audience profiles.

\subsection{Runtime Profile}

The runtime profile is summarized in Table~\ref{tab:significance_cost}. EAI-CO required 8.0 mean model calls per task and 11114.23 ms mean latency, whereas the single-pass baselines required approximately 3--4 model calls and 2700--2800 ms mean latency. This overhead is expected because the framework explicitly trades additional inference budget for exploratory search quality. For the intended offline campaign-design setting, this tradeoff remains operationally acceptable.

\subsection{Cross-model Validation on \texttt{gpt-5.4-mini-2026-03-17}}

Table~\ref{tab:gpt_cross_model} presents the fixed-subset cross-model validation on \texttt{gpt-5.4-mini-2026-03-17}. The same ranking pattern observed on the local backbone was preserved: EAI-CO achieved the highest reward (\textbf{0.9957}), followed by \texttt{B2\_OpenSource\_Only} (0.9751), \texttt{B1\_SingleShot\_API} (0.9736), \texttt{B3\_PromptEngineered\_AI} (0.9729), and \texttt{B0\_Template} (0.8396). The gains over all baselines remained statistically significant, with reward improvements of +0.1561 ($p=1.54\times10^{-8}$) over \texttt{B0\_Template}, +0.0221 ($p=1.41\times10^{-3}$) over \texttt{B1\_SingleShot\_API}, +0.0206 ($p=6.42\times10^{-4}$) over \texttt{B2\_OpenSource\_Only}, and +0.0228 ($p=1.01\times10^{-3}$) over \texttt{B3\_PromptEngineered\_AI}.

The audience-wise rankings were also preserved, with EAI-CO remaining first for family users (0.9950), price-sensitive consumers (0.9938), students (1.0000), and young professionals (0.9939). These results support the claim that the advantage of the proposed exploratory optimization framework transfers to a stronger commercial model under the same task protocol.

\section{Discussion}

\subsection{The Main Advantage Comes from Optimization Structure}

The most stable pattern across the local benchmark and the commercial-model validation is that exploratory optimization outperforms one-shot generation. This result suggests that the central contribution of EAI-CO lies less in any specific generator and more in the optimization structure itself. In digital-media advertising, creative quality should therefore be treated as a search-and-selection problem rather than as a direct text-generation problem.

\subsection{Why Diversity and Iteration Matter}

The ablation study shows that diversity preservation and iterative refinement are the main structural sources of improvement. This is consistent with how advertising work is performed in practice. Strong campaign creatives rarely emerge from the first acceptable draft; they emerge from comparing strategically different variants and refining the strongest direction. The present benchmark makes this process measurable.

\subsection{Audience Modeling Has Measurable Functional Value}

The audience-modeling ablation provides direct empirical evidence that personalization is not simply a rhetorical feature. Once persona conditioning is removed, both reward and audience fit decline. This implies that the framework is not merely generating generic persuasive copy but is measurably exploiting audience-specific information to alter the optimization trajectory.

\subsection{Compliance-aware Scoring Is a Practical Requirement}

The factual-penalty ablation remained close to the full framework, which indicates that compliance-aware optimization can be introduced without substantially damaging persuasive performance. This is a practically desirable result. In real digital-media settings, a method that produces slightly stronger promotional language at the cost of weaker factual discipline would not necessarily be preferable.

\subsection{Cross-model Transfer Strengthens the Paper}

The commercial-model experiment should be interpreted as a transferability test rather than as the main source of effect-size contrast. Because \texttt{gpt-5.4-mini} is a stronger generator under the current evaluator, reward values become compressed near the upper bound. Nevertheless, the preservation of the full ranking pattern materially strengthens the paper by reducing the risk that the main result is merely an artifact of one local open-source backbone.

\subsection{Limitations}

The present study does not include real platform CTR or conversion outcomes, and the current submission package does not yet include a completed human-evaluation study. Accordingly, the paper should restrict its claims to automatic reward superiority, predicted engagement, audience-conditioned quality, and cross-model ranking consistency. In addition, the current benchmark remains smaller than a production-scale industrial evaluation environment, and future work should expand both product diversity and deployment realism.


\section{Results}

\subsection{Main Benchmark Results on the Local Backbone}

Table~\ref{tab:qwen_main_results} reports the main benchmark results obtained on the local \texttt{Qwen2.5-7B-Instruct} backbone. EAI-CO achieved the highest overall reward (\textbf{0.8051}), outperforming \texttt{B3\_PromptEngineered\_AI} (0.7568), \texttt{B1\_SingleShot\_API} (0.7494), \texttt{B2\_OpenSource\_Only} (0.7445), and \texttt{B0\_Template} (0.6856). As summarized in Table~\ref{tab:significance_cost}, the reward gains over all primary baselines were statistically significant. Relative to \texttt{B0\_Template}, EAI-CO improved reward by +0.1195 ($p=9.09\times10^{-13}$). The gains over the stronger single-pass baselines also remained significant: +0.0557 over \texttt{B1\_SingleShot\_API} ($p=1.34\times10^{-6}$), +0.0606 over \texttt{B2\_OpenSource\_Only} ($p=5.18\times10^{-8}$), and +0.0483 over \texttt{B3\_PromptEngineered\_AI} ($p=6.58\times10^{-5}$).

Beyond the scalar reward, the full framework achieved the best audience-fit score (0.4719) and the strongest brand-safety score (0.9945) among the primary methods. This indicates that the observed advantage is not restricted to a single dimension such as relevance or engagement, but reflects a more balanced optimization outcome.

\subsection{Ablation Analysis}

Table~\ref{tab:qwen_ablation} presents the ablation study on the same benchmark. The full method outperformed all ablated variants. The largest degradation occurred when diversity preservation was removed (-0.0702), indicating that search-space breadth is a major contributor to the final gain. Removing the iterative loop reduced reward by -0.0291, showing that controlled refinement contributes beyond one-pass generation. Removing audience modeling lowered reward by -0.0218 and reduced audience fit from 0.4719 to 0.3960, directly supporting the claim that persona conditioning is functionally important.

The factual-penalty ablation remained close to the full model (-0.0045). This result should be interpreted carefully. It does not imply that compliance-aware controls are unnecessary. Rather, it indicates that factual safeguards can be integrated with only a small tradeoff in persuasive reward while preserving a more disciplined search process.

\subsection{Audience-wise Robustness}

Table~\ref{tab:qwen_audience} reports the audience-wise reward values on the local benchmark. EAI-CO ranked first in all four audience groups: family users (0.8147), students (0.8098), young professionals (0.8050), and price-sensitive consumers (0.7908). This consistency reduces the concern that the overall effect is driven by a single favorable subgroup and suggests that the exploratory optimization strategy remains effective across materially different audience profiles.

\subsection{Runtime Profile}

The runtime profile is summarized in Table~\ref{tab:significance_cost}. EAI-CO required 8.0 mean model calls per task and 11114.23 ms mean latency, whereas the single-pass baselines required approximately 3--4 model calls and 2700--2800 ms mean latency. This overhead is expected because the framework explicitly trades additional inference budget for exploratory search quality. For the intended offline campaign-design setting, this tradeoff remains operationally acceptable.

\subsection{Cross-model Validation on \texttt{gpt-5.4-mini-2026-03-17}}

Table~\ref{tab:gpt_cross_model} presents the fixed-subset cross-model validation on \texttt{gpt-5.4-mini-2026-03-17}. The same ranking pattern observed on the local backbone was preserved: EAI-CO achieved the highest reward (\textbf{0.9957}), followed by \texttt{B2\_OpenSource\_Only} (0.9751), \texttt{B1\_SingleShot\_API} (0.9736), \texttt{B3\_PromptEngineered\_AI} (0.9729), and \texttt{B0\_Template} (0.8396). The gains over all baselines remained statistically significant, with reward improvements of +0.1561 ($p=1.54\times10^{-8}$) over \texttt{B0\_Template}, +0.0221 ($p=1.41\times10^{-3}$) over \texttt{B1\_SingleShot\_API}, +0.0206 ($p=6.42\times10^{-4}$) over \texttt{B2\_OpenSource\_Only}, and +0.0228 ($p=1.01\times10^{-3}$) over \texttt{B3\_PromptEngineered\_AI}.

The audience-wise rankings were also preserved, with EAI-CO remaining first for family users (0.9950), price-sensitive consumers (0.9938), students (1.0000), and young professionals (0.9939). These results support the claim that the advantage of the proposed exploratory optimization framework transfers to a stronger commercial model under the same task protocol.

\section{Discussion}

\subsection{The Main Advantage Comes from Optimization Structure}

The most stable pattern across the local benchmark and the commercial-model validation is that exploratory optimization outperforms one-shot generation. This result suggests that the central contribution of EAI-CO lies less in any specific generator and more in the optimization structure itself. In digital-media advertising, creative quality should therefore be treated as a search-and-selection problem rather than as a direct text-generation problem.

\subsection{Why Diversity and Iteration Matter}

The ablation study shows that diversity preservation and iterative refinement are the main structural sources of improvement. This is consistent with how advertising work is performed in practice. Strong campaign creatives rarely emerge from the first acceptable draft; they emerge from comparing strategically different variants and refining the strongest direction. The present benchmark makes this process measurable.

\subsection{Audience Modeling Has Measurable Functional Value}

The audience-modeling ablation provides direct empirical evidence that personalization is not simply a rhetorical feature. Once persona conditioning is removed, both reward and audience fit decline. This implies that the framework is not merely generating generic persuasive copy but is measurably exploiting audience-specific information to alter the optimization trajectory.

\subsection{Compliance-aware Scoring Is a Practical Requirement}

The factual-penalty ablation remained close to the full framework, which indicates that compliance-aware optimization can be introduced without substantially damaging persuasive performance. This is a practically desirable result. In real digital-media settings, a method that produces slightly stronger promotional language at the cost of weaker factual discipline would not necessarily be preferable.

\subsection{Cross-model Transfer Strengthens the Paper}

The commercial-model experiment should be interpreted as a transferability test rather than as the main source of effect-size contrast. Because \texttt{gpt-5.4-mini} is a stronger generator under the current evaluator, reward values become compressed near the upper bound. Nevertheless, the preservation of the full ranking pattern materially strengthens the paper by reducing the risk that the main result is merely an artifact of one local open-source backbone.

\subsection{Limitations}

The present study does not include real platform CTR or conversion outcomes, and the current submission package does not yet include a completed human-evaluation study. Accordingly, the paper should restrict its claims to automatic reward superiority, predicted engagement, audience-conditioned quality, and cross-model ranking consistency. In addition, the current benchmark remains smaller than a production-scale industrial evaluation environment, and future work should expand both product diversity and deployment realism.
